% Requires: \usepackage{booktabs}
% All values below are from the PartImageNet++ 10k annotation holdout.

\subsection{PartImageNet++ Results and Evaluation Caveat}

All methods use the same scratch ResNet-50 backbone trained on the 90k
PartImageNet++ training split and the same 783 generic concepts.  NEC uses a
SAGA/GLM sparse linear classifier fit on the concept activations.  Concept
presence is evaluated against human COCO boxes.  The localization evaluation
uses human polygons or boxes after the same \texttt{Resize(256)} and
\texttt{CenterCrop(224)} transform applied to the models.  We exclude 1,120
identity targets (4.19\% of the original image--concept targets), such as the
whole-object label \emph{bakery} for a bakery image.

\begin{table*}[t]
\centering
\caption{Classification and concept-prediction results on the 10k PIN++
annotation holdout.  NEC columns are top-1 accuracy (\%).  Concept-presence
metrics use filtered human GT boxes; higher is better for every metric.}
\label{tab:pinpp-classification-concepts}
\small
\setlength{\tabcolsep}{5pt}
\begin{tabular}{lrrrrrr|rrrrr}
\toprule
& \multicolumn{6}{c|}{NEC top-1 (\%)}
& \multicolumn{5}{c}{Human-box concept presence} \\
\cmidrule(lr){2-7}\cmidrule(l){8-12}
Method & 5 & 10 & 15 & 20 & 25 & 30
& AUROC & AP & Macro AP & P@5 & Best F1 \\
\midrule
SALF-CBM & 23.41 & 40.99 & 48.37 & 52.08 & 54.43 & 55.41
& 0.5790 & 0.0047 & 0.0482 & 0.0063 & 0.0114 \\
VLG-CBM & 19.05 & 24.56 & 29.36 & 32.44 & 35.40 & 37.49
& \textbf{0.9945} & \textbf{0.6190} & \textbf{0.5016} & \textbf{0.4144} & \textbf{0.6555} \\
SG-CBM & \textbf{47.02} & \textbf{55.08} & \textbf{58.29} & \textbf{59.34} & \textbf{59.81} & \textbf{60.21}
& \textbf{0.9945} & 0.6075 & 0.4861 & 0.4055 & 0.6324 \\
\bottomrule
\end{tabular}
\end{table*}

\begin{table*}[t]
\centering
\caption{Human localization results on the same filtered 10k holdout.  Polygon
metrics use human segmentation masks; box metrics use human COCO boxes.
``Box IoU'' is the IoU of the tight predicted activation box and GT box at the
per-map mean threshold.  Raw RMA is marked with $\dagger$ because spatial
softmax mass depends on the arbitrary scale of each method's map logits.}
\label{tab:pinpp-localization}
\small
\setlength{\tabcolsep}{5pt}
\begin{tabular}{lrrrrrr|rrrrr}
\toprule
& \multicolumn{6}{c|}{Human polygons at $224\times224$}
& \multicolumn{5}{c}{Human boxes at $224\times224$} \\
\cmidrule(lr){2-7}\cmidrule(l){8-12}
Method & RMA$^\dagger$ & Pointing & Soft IoU & mIoU & Dice & Pixel mAP
& RMA$^\dagger$ & Pointing & Soft IoU & Box IoU & Acc@0.5 \\
\midrule
SALF-CBM & 0.1583 & 0.2532 & 0.1051 & 0.1790 & 0.2781 & 0.2501
& 0.2515 & 0.3642 & 0.1801 & 0.2921 & 0.2243 \\
VLG-CBM & \textbf{0.2811} & 0.2999 & 0.1004 & 0.1848 & 0.2851 & 0.2695
& \textbf{0.4195} & 0.4412 & 0.1282 & 0.3051 & 0.2406 \\
SG-CBM & 0.2749 & \textbf{0.4117} & \textbf{0.1774} & \textbf{0.2177} & \textbf{0.3231} & \textbf{0.3153}
& 0.4097 & \textbf{0.5681} & \textbf{0.2629} & \textbf{0.3239} & \textbf{0.2693} \\
\bottomrule
\end{tabular}
\end{table*}

\paragraph{Interpretation and limitation.}
PIN++ provides high-quality, class-specific human part annotations, but our
derived 783-concept benchmark collapses the source vocabulary of 3,310
class-specific part categories into generic names.  This loses the relation
``\emph{part of which object}?''  For example, on an image containing a killer
whale and people, a generic \emph{head} map that activates on a person's head
is semantically reasonable but is scored as false because the annotation only
labels the killer whale's head.  Thus the current unconditioned generic-part
protocol conflates target-object association with generic-part localization.

The generic vocabulary is also incomplete for this use case: some source
categories deliberately use the entire foreground object as one ``part'', and
target size varies substantially after preprocessing (13.8\% of targets cover
less than 1\% of the crop, whereas 8.1\% cover more than half).  The identity
filter removes the clearest whole-object cases but cannot resolve generic-part
ambiguity or distractor instances.  Finally, this 10k split is a 9:1 holdout
within the PIN++ annotations drawn from ImageNet training images, not the
official ImageNet validation set.

\paragraph{Recommended positioning.}
These experiments are useful controlled evidence that SG-CBM has substantially
better sparse classification and spatial overlap than the two baselines under
the current protocol.  They should not be the sole evidence for a generic
localization claim.  A stronger primary evaluation should be class-conditioned
(e.g., \emph{killer whale head}) or explicitly target-object-conditioned, with
an independently annotated ImageNet-validation test set, target-size strata,
and a released manual QA audit of the genericization mapping.
